\section{Analysis}
\label{sec:analysis}

The preceding experiments reveal a sharp divergence: pruning any fraction of SmolVLA's vision tokens causes complete task failure, while OpenVLA degrades gracefully under the same operations, and learned compression on SmolVLA even exceeds the uncompressed baseline. This section traces these outcomes to three architectural factors: the connector design that determines token granularity (\S\ref{sec:pixel_shuffle}), the attention distribution that reflects downstream reliance on individual tokens (\S\ref{sec:attention_distribution}), and the mechanism by which learned compression sidesteps the failure mode of hard pruning (\S\ref{sec:learned_compression}).

\subsection{The Pixel-Shuffle Bottleneck}
\label{sec:pixel_shuffle}

SmolVLA's vision encoder (SigLIP) produces 1024 patches per camera image on a $32 \times 32$ grid. Before these patches reach the language backbone, a pixel-shuffle connector applies $4 \times 4$ spatial folding, mapping every $4 \times 4$ block of patches into a single token by channel concatenation. Each camera's 1024 patches are thereby compressed into 64 tokens, and the dual-camera setup yields 128 vision tokens in total. The critical consequence is that each token after pixel-shuffle encodes the full semantic content of 16 spatially contiguous patches. Removing a single token does not discard a localized detail; it eliminates an entire $4 \times 4$ receptive field with no redundancy to compensate for the loss. Because the folding is a fixed, non-learned operation applied before the language model, no downstream layer can recover the discarded information. This creates a cliff-edge failure mode: at every pruning ratio we tested, SmolVLA's success rate drops to 0\%.

OpenVLA uses a qualitatively different connector. Its MLP projector applies an independent linear transformation to each SigLIP patch, producing 256 vision tokens that preserve the one-to-one correspondence between tokens and spatial locations. Removing one token discards approximately $1/256$ of the visual field, a localized loss that neighboring tokens can partially compensate for through the self-attention layers that follow. At 50\% pruning ($k=0.5$), OpenVLA retains 37.4\% success rate against a 74\% baseline, a substantial but not catastrophic reduction. The contrast is architectural, not a matter of model scale: pixel-shuffle multiplies the per-token information density by $16\times$, and this multiplicative factor converts any token-level removal into a large, irreversible information gap. While model scale is a confound (0.5B vs 7B), the extreme sensitivity of SmolVLA, collapsing even at 90\% token retention ($k=0.9$, removing only 13 of 128 tokens), is too severe to explain by capacity alone.

\subsection{Attention Distribution as a Compressibility Predictor}
\label{sec:attention_distribution}

The attention entropy ratio provides a quantitative lens on how the action decoder distributes its reliance across vision tokens. We define the entropy ratio as the normalized entropy of the attention weights that the decoder's action head assigns to vision token positions, where 1.0 corresponds to a perfectly uniform distribution and lower values indicate selective concentration.

SmolVLA exhibits an overall entropy ratio of 0.952 ($\pm$0.012), measured on 260 real LIBERO-Object observations across all 10 tasks. The per-layer profile is remarkably flat: entropy ratios range from 0.931 (layer 13) to 0.983 (layer 0), a total variation of only 0.05 across all 16 transformer layers. No layer develops selective attention; the action decoder assigns near-equal weight to all 128 vision tokens at every depth. This observation is consistent with the pixel-shuffle analysis above. Because each token already encodes 16 spatially contiguous patches, every token is information-dense, and no downstream layer learns to treat any token as dispensable. Under these conditions, pruning any subset of tokens removes a proportional fraction of the decoder's input signal, and the 0\% success rate across all pruning methods and ratios follows directly.

OpenVLA tells a different story. Measured on 250 real LIBERO-Object observations spanning all 10 tasks, its overall entropy ratio is 0.639 ($\pm$0.133), and the layer-wise breakdown reveals progressive sharpening: shallow layers (0\textendash 1) show ratios of 0.98\textendash 0.95, mid layers (5\textendash 16) drop to 0.62\textendash 0.69, and the deepest layers (27\textendash 30) drop further to 0.48\textendash 0.54, with layer 29 reaching a minimum ratio of 0.477. The total variation across OpenVLA's 32 layers is 0.50, an order of magnitude larger than SmolVLA's 0.05. These deep layers, closest to the action prediction head, concentrate the majority of their attention on a small subset of the 256 vision tokens. Only 12.1\% of total attention weight falls on vision tokens overall, and within that fraction, the distribution is highly non-uniform. This pattern indicates that OpenVLA has learned to identify which spatial locations carry task-relevant information and to downweight the rest. Pruning low-attention tokens therefore removes positions that the model already ignores, and the resulting performance loss is proportional to the actual information discarded rather than to the count of tokens removed.

These findings suggest that the attention entropy ratio can serve as a practical, preliminary diagnostic signal for compressibility. A model with entropy ratio close to 1.0 distributes reliance uniformly and will not tolerate token removal; a model with substantially lower entropy ratio has internalized a sparse selection mechanism and can absorb pruning with graceful degradation. This metric can be computed from a single forward pass on a small calibration set, making it a lightweight diagnostic before committing to a compression strategy.

\subsection{Why Learned Compression Bypasses the Cliff-Edge}
\label{sec:learned_compression}

If hard pruning fails on SmolVLA because every token is indispensable, why does learned compression with cross-attention pooling succeed? The answer lies in the distinction between deletion and aggregation. Hard pruning applies a binary mask: selected tokens are kept verbatim, and the rest are discarded with their information permanently lost. Cross-attention pooling operates differently. A set of $M$ learned query vectors attends over all 128 source tokens through standard cross-attention, and each query computes a weighted sum across the full token set. Even when the output is compressed to $M = 32$ (a 75\% reduction in token count), every original token contributes to the compressed representation through its attention weight. No information is deleted; it is redistributed into fewer, higher-capacity vectors.

The experimental results confirm this mechanism across five settings, including a no-compression control. With $M = 32$ (75\% token reduction), Smart Compress achieves 61.6\% SR (308/500, CI [57.3\%, 65.8\%]), and $M = 64$ (50\% reduction) reaches the same 61.6\% SR (308/500, CI [57.3\%, 65.8\%]). At more extreme compression, $M = 16$ (87.5\% reduction) achieves 55.8\% SR (279/500, CI [51.4\%, 60.1\%]), while $M = 96$ (25\% reduction) yields 52.6\% SR (263/500, CI [48.2\%, 56.9\%]). The passthrough control $M = 128$ (no compression, same token count as the original input) achieves 59.8\% SR (299/500, CI [55.4\%, 64.0\%]). All five configurations exceed the uncompressed baseline of 44.8\% by 7.8 to 16.8 percentage points, in the same setting where hard pruning produces 0\% SR at every ratio tested.

The $M = 128$ result is the most informative control. Because $M = 128$ equals the original token count, the cross-attention module performs no compression; it only re-encodes the tokens through learned queries. The 15.0 percentage point improvement over baseline (59.8\% vs 44.8\%) demonstrates that the cross-attention architecture itself is the primary source of benefit, acting as a learned adapter that re-encodes the pixel-shuffle output into a more action-relevant representation. The compressed settings $M = 32$ and $M = 64$ achieve a further 1.8 percentage points above $M = 128$ (61.6\% vs 59.8\%), though the confidence intervals overlap, suggesting that the compression bottleneck provides at most a modest additional regularization effect. The per-task distributions differ across all five settings (Table~\ref{tab:smart_compress}), indicating that each compression ratio allocates information differently across tasks.

Two factors explain why the cross-attention module exceeds the baseline. First, the compressor (22.2M parameters) introduces representational capacity that the frozen pixel-shuffle connector lacks; it learns a task-adapted re-encoding of the visual input during training. Second, at moderate compression ratios ($M = 32$--$64$), the bottleneck additionally filters noise and redundancy from the vision tokens, retaining task-relevant features while suppressing distractors. The combination of learned re-encoding (present at all $M$ values) and information compression (present when $M < 128$) converts the same 128 tokens that resist pruning into a more useful representation.

This result carries a practical implication for VLA deployment. Architectures with aggressive spatial folding in their vision connectors are not inherently incompressible; they are incompressible under deletion-based methods. A learned compression module that aggregates rather than prunes can recover and even exceed baseline performance, offering a viable path to reduced inference cost without the cliff-edge failure that hard pruning induces.
